\documentclass[a4paper,12pt]{article}
\usepackage[T2A]{fontenc}
\usepackage[utf8]{inputenc}
\usepackage[russian]{babel}
\usepackage{amsmath,amssymb,icomma}
\usepackage[margin=20mm,top=23mm,bottom=22mm,headheight=26pt]{geometry}
\usepackage{microtype,parskip,setspace,fancyhdr,booktabs,tabularx,enumitem,xcolor,titlesec}
\usepackage{tikz,tkz-euclide}
\usetikzlibrary{arrows.meta,positioning,shapes.geometric}
\usepackage[hidelinks]{hyperref}
\definecolor{accent}{HTML}{008080}
\definecolor{darkaccent}{HTML}{004D4D}
\definecolor{textgray}{HTML}{333333}
\color{textgray}
\onehalfspacing
\setlength{\parskip}{5pt}
\setlength{\parindent}{0pt}
\setlist{nosep,leftmargin=*}
\titleformat{\section}{\Large\color{darkaccent}\normalfont}{\thesection.}{0.5em}{}[\vspace{3pt}\color{accent}\titlerule]
\titlespacing*{\section}{0pt}{0pt}{9pt}
\newcommand{\sub}[1]{\par\medskip\textcolor{darkaccent}{\textbf{#1}}\par\smallskip}
\newcommand{\Sos}{S_{\text{осн}}}
\newcommand{\Sb}{S_{\text{бок}}}
\newcommand{\Sp}{S_{\text{полн}}}
\pagestyle{fancy}
\fancyhf{}
\fancyhead[L]{\small Пирамиды, конус и составные тела}
\fancyhead[R]{\small Николь Саркисьянц}
\fancyfoot[C]{\small\thepage}
\renewcommand{\headrulewidth}{0.4pt}
\renewcommand{\headrule}{\hbox to\headwidth{\color{accent}\leaders\hrule height \headrulewidth\hfill}}
\begin{document}
\begin{titlepage}
\centering
\vspace*{26mm}
{\Large\color{darkaccent} МАТЕМАТИКА\par}
\vspace{14mm}
{\Huge Чек-лист\par}
\vspace{10mm}
{\LARGE\color{accent} Пирамиды, конус\par и составные тела\par}
\vspace{13mm}
{\large Подготовка к ЕГЭ по профильной математике\par}
\vfill
{\large Николь Саркисьянц\par}
\vspace{10mm}
{26 ноября 2026 г.\par}
\vspace{16mm}
{\small Лёвин Артём Александрович\par эксклюзивно для Николь Саркисьянц\par}
\vspace{17mm}
\begin{tikzpicture}
\draw[accent,line width=.7pt] (0,0) .. controls (2,0.8) and (4,-0.8) .. (6,0) .. controls (8,0.8) and (10,-0.8) .. (12,0);
\end{tikzpicture}
\end{titlepage}
\setcounter{page}{1}
\section{Правильная четырёхугольная пирамида}
В основании правильной четырёхугольной пирамиды лежит квадрат. Вершина $S$ проецируется в его центр $O$. Обозначим сторону основания через $a$, диагональ через $d$, середину стороны через $M$.
\[
\Sos=a^2=\frac{d^2}{2},\qquad d=a\sqrt2,\qquad P_{\text{осн}}=4a.
\]
\begin{center}
\begin{tikzpicture}[scale=.83,line width=.65pt]
\tkzDefPoint(0,0){A}\tkzDefPoint(5,0){B}\tkzDefPoint(6.8,2){C}\tkzDefPoint(1.8,2){D}
\tkzDefMidPoint(A,C)\tkzGetPoint{O}
\tkzDefMidPoint(A,B)\tkzGetPoint{M}
\tkzDefPoint(3.4,4.7){S}
\tkzDrawSegments(A,B B,C C,S S,A S,B S,M)
\tkzDrawSegments[dashed,gray](A,D D,C S,D A,C B,D)
\tkzDrawSegments[color=accent,line width=1pt](S,O O,M)
\tkzDrawPoints(A,B,C,D,O,M,S)
\tkzLabelPoints[below](A,B,M)\tkzLabelPoints[right](C,O)\tkzLabelPoints[left](D)\tkzLabelPoints[above](S)
\end{tikzpicture}
\hspace{8mm}
\begin{tikzpicture}[scale=.83,line width=.7pt]
\tkzDefPoint(0,0){O}\tkzDefPoint(2.8,0){M}\tkzDefPoint(0,2.8){S}
\tkzDrawPolygon(S,O,M)\tkzMarkRightAngle[size=.25](S,O,M)
\tkzMarkAngle[size=.55,color=accent](S,M,O)\tkzLabelAngle[pos=.85](S,M,O){$\alpha$}
\tkzLabelPoints[below](O,M)\tkzLabelPoints[above](S)
\node[left] at (0,1.4){$H$};\node[above right] at (1.4,1.4){$l$};\node[below] at (1.4,-.35){$a/2$};
\end{tikzpicture}
\end{center}
\smallskip
\begin{tabularx}{\linewidth}{@{}lX@{}}
\toprule
Отрезок & Смысл \\\midrule
$SO=H$ & Высота пирамиды, перпендикулярная основанию.\\
$SM=l$ & Апофема, высота боковой грани.\\
$SA=b$ & Боковое ребро, соединяющее вершину с вершиной основания.\\
$OA=d/2$ & Половина диагонали квадрата.\\
$OM=a/2$ & Расстояние от центра квадрата до стороны.\\\bottomrule
\end{tabularx}
\[
V=\frac13\Sos H,\qquad \Sb=\frac12P_{\text{осн}}l,\qquad \Sp=\Sb+\Sos.
\]
\textbf{Проверяйте обозначения:} для объёма используется $H$; для боковой поверхности используется $l$. Эти отрезки имеют разную длину.
\clearpage
\section{Как найти высоту и апофему}
\sub{Пример 1. Дана диагональ основания и боковое ребро}
Диагональ квадратного основания равна $10$, боковое ребро равно $\sqrt{34}$. Найдите объём.
\begin{center}
\begin{tikzpicture}[scale=.65]
\tkzDefPoint(0,0){O}\tkzDefPoint(5,0){A}\tkzDefPoint(0,3){S}
\tkzDrawPolygon[color=darkaccent](S,O,A)\tkzMarkRightAngle[size=.35](S,O,A)
\tkzLabelPoints[below](O,A)\tkzLabelPoints[above](S)
\node[left] at (0,1.5){$H$};\node[below] at (2.5,0){$5$};\node[above right] at (2.5,1.5){$\sqrt{34}$};
\end{tikzpicture}
\end{center}
\[
\Sos=\frac{10^2}{2}=50,\qquad OA=\frac{10}{2}=5.
\]
В прямоугольном треугольнике $SOA$ гипотенуза $SA$, поэтому
\[
H=\sqrt{SA^2-OA^2}=\sqrt{34-25}=3,\quad V=\frac13\cdot50\cdot3=50.
\]
\sub{Пример 2. Дана сторона основания и двугранный угол}
Сторона основания равна $12$, угол между боковой гранью и плоскостью основания равен $45^\circ$. Найдите объём и полную площадь поверхности.

$SM$ и $OM$ перпендикулярны стороне основания, поэтому $\angle SMO=\alpha$ --- линейный угол двугранного угла. В треугольнике $SOM$:
\[
\tg\alpha=\frac{SO}{OM},\qquad \cos\alpha=\frac{OM}{SM}.
\]
Следовательно,
\[
H=6\tg45^\circ=6,\qquad l=\frac{6}{\cos45^\circ}=6\sqrt2.
\]
\[
V=\frac13\cdot144\cdot6=288,
\]
\[
\Sp=\frac12\cdot48\cdot6\sqrt2+144=144(\sqrt2+1).
\]
\textbf{Самопроверка:} $l^2=H^2+(a/2)^2$, следовательно $l>H$.
\clearpage
\section{Тетраэдр и изменение размеров}
Правильный тетраэдр --- пирамида, все шесть рёбер которой равны. Каждая грань является равносторонним треугольником со стороной $a$.
\[
\Sos=\frac{\sqrt3}{4}a^2,\qquad \Sp=4\Sos=\sqrt3 a^2.
\]
Расстояние от центра основания до его вершины равно $a/\sqrt3$. По теореме Пифагора высота тетраэдра равна
\[
H=\sqrt{a^2-\frac{a^2}{3}}=a\sqrt{\frac23}.
\]
Подставляем в формулу объёма:
\[
V=\frac13\cdot\frac{\sqrt3}{4}a^2\cdot a\sqrt{\frac23}
=\frac{\sqrt2}{12}a^3.
\]
\sub{Подобные тела}
Если все соответствующие линейные размеры умножаются на $k$ ($k>0$), то
\[
\frac{S_{\text{нов}}}{S_{\text{исх}}}=k^2,\qquad
\frac{V_{\text{нов}}}{V_{\text{исх}}}=k^3.
\]
Например, при удвоении ребра правильного тетраэдра площадь увеличится в $2^2=4$ раза, объём --- в $2^3=8$ раз. При утроении ребра объём увеличится в $27$ раз.
\sub{Независимое изменение основания и высоты}
Для пирамиды с квадратным основанием $V=\frac13 a^2H$. Если сторона основания увеличилась в $p$ раз, а высота в $q$ раз, то
\[
\frac{V_{\text{нов}}}{V_{\text{исх}}}=\frac{(pa)^2qH}{a^2H}=p^2q.
\]
\begin{center}
\begin{tabular}{@{}lll@{}}
\toprule
Сторона основания & Высота & Объём \\\midrule
в $2$ раза & в $3$ раза & в $2^2\cdot3=12$ раз\\
в $3$ раза & в $2$ раза & в $3^2\cdot2=18$ раз\\\bottomrule
\end{tabular}
\end{center}
\textbf{Важно:} если $p\ne q$, подобие тел в общем случае отсутствует. Для полной площади поверхности тогда нельзя автоматически применять множитель $p^2$ или $q^2$.
\clearpage
\section{Составные тела: вырезанный угол}
Для прямоугольного параллелепипеда с измерениями $a,b,c$:
\[
V_0=abc,\qquad S_0=2(ab+ac+bc).
\]
При удалении части тела вычитаем её объём. При присоединении части без объёмного перекрытия прибавляем её объём.
\[
V=V_0-V_{\text{выр}}+V_{\text{добав}}.
\]
Для площади поверхности удобен отдельный баланс:
\[
\Sp=S_0-S_{-}+S_{+}.
\]
Здесь $S_-$ --- участки прежней поверхности, исчезнувшие после преобразования; $S_+$ --- новые открытые участки. Внутренние грани соприкасающихся деталей не входят во внешнюю поверхность.
\sub{Пример 3. Из угла вырезан куб}
Из угла параллелепипеда $4\times3\times2$ вырезали куб с ребром $1$. Удалены участки трёх внешних граней; образованы три новые квадратные площадки той же площади.
\begin{center}
\begin{tikzpicture}[scale=.85]
\tkzDefPoint(0,0){A}\tkzDefPoint(4,0){B}\tkzDefPoint(5.4,1){C}\tkzDefPoint(0,2){D}\tkzDefPoint(4,2){E}\tkzDefPoint(5.4,3){F}\tkzDefPoint(1.4,3){G}
\tkzDefPoint(0,1){I}\tkzDefPoint(1,1){J}\tkzDefPoint(1,2){K}\tkzDefPoint(.467,2.333){L}\tkzDefPoint(1.467,2.333){N}\tkzDefPoint(.467,1.333){P}\tkzDefPoint(1.467,1.333){Q}
\tkzDrawPolygon[fill=gray!7](A,B,E,K,J,I)
\tkzDrawPolygon[fill=gray!12](B,C,F,E)
\tkzDrawPolygon[fill=gray!5](L,G,F,E,K,N)
\tkzDrawPolygon[fill=accent!15](I,J,Q,P)
\tkzDrawPolygon[fill=accent!27](J,K,N,Q)
\tkzDrawPolygon[fill=accent!10](P,Q,N,L)
\tkzDrawSegments(A,I I,P P,L)
\node at (3.25,1){$4\times3\times2$};
\node[accent,above] at (1.2,3.1){вырез};
\draw[accent,-{Stealth}] (1.2,3.05)--(1.1,1.95);
\end{tikzpicture}
\end{center}
\[
V=4\cdot3\cdot2-1^3=23.
\]
\[
\Sp=2(4\cdot3+4\cdot2+3\cdot2)-3\cdot1^2+3\cdot1^2=52.
\]
В общем случае для углового выреза куба с ребром $t$, если $0<t<\min(a,b,c)$,
\[
V=abc-t^3,\qquad \Sp=2(ab+ac+bc).
\]
\textbf{Идея урока:} мысленно перенесите три исчезнувшие площадки на стенки выреза. Их суммарная площадь сохраняется.
\clearpage
\section{Сквозной паз и присоединённый куб}
Пусть в параллелепипеде $a\times b\times c$ выполнен прямоугольный паз длиной $a$, шириной $w$ и глубиной $h$. Паз открыт сверху и с двух торцов, расположен строго между боковыми гранями; $0<w<b$, $0<h<c$.
\begin{center}
\begin{tikzpicture}[scale=.8]
\tkzDefPoint(0,0){A}\tkzDefPoint(6,0){B}\tkzDefPoint(6,3){C}\tkzDefPoint(4,3){D}\tkzDefPoint(4,1){E}\tkzDefPoint(2,1){F}\tkzDefPoint(2,3){G}\tkzDefPoint(0,3){H}
\tkzDrawPolygon[fill=gray!8](A,B,C,D,E,F,G,H)
\tkzDrawSegments[color=accent,line width=1.3pt](D,E E,F F,G)
\node[above] at (3,3.05){$w$};\draw[<->,accent] (2,3.02)--(4,3.02);
\draw[<->,accent] (3,1.1)--(3,2.9);\node[right] at (3,2){$h$};
\node[below] at (3,0){$b$};\node[left] at (0,1.5){$c$};
\end{tikzpicture}
\par\small Поперечное сечение, перпендикулярное длине паза.
\end{center}
\begin{tabularx}{\linewidth}{@{}Xr@{}}
\toprule
Изменение поверхности & Площадь\\\midrule
Исчезли два участка торцов & $-2wh$\\
Появились две продольные стенки & $+2ah$\\
Полоса сверху заменена дном паза & $-aw+aw=0$\\\bottomrule
\end{tabularx}
\[
V=abc-awh,\qquad \Sp=2(ab+ac+bc)+2h(a-w).
\]
\sub{Пример 4. Параллелепипед $5\times4\times3$, паз $5\times1\times2$}
\[
V=60-10=50,\quad \Sp=94-4+20=110.
\]
\sub{Пример 5. Паз и куб на верхней грани}
Параллелепипед $6\times4\times3$, паз $6\times1\times2$, куб с ребром $1$ полностью прилегает нижней гранью к свободной части верха и не перекрывает паз.
\[
V=72-12+1=61,\quad \Sp=108+2\cdot2(6-1)+4=132.
\]
При присоединении куба с ребром $t$ исчезает площадка $t^2$ и появляются пять его граней, поэтому прирост площади равен $5t^2-t^2=4t^2$.
\clearpage
\section*{Алгоритм: объём и площадь составного тела}
\vspace{9mm}
\begin{center}
\begin{tikzpicture}[node distance=9mm,>=Stealth,
start/.style={ellipse,draw=accent,align=center,minimum width=42mm,minimum height=9mm},
actionbox/.style={rectangle,rounded corners=2pt,draw=accent,align=center,text width=62mm,minimum height=12mm,inner sep=5pt},
choice/.style={diamond,aspect=2.4,draw=accent,align=center,text width=37mm,inner sep=1pt},
arr/.style={-{Stealth[length=2.7mm]},draw=darkaccent,line width=.6pt}]
\node[start] (start) {Начало};
\node[actionbox,below=of start] (base) {Найдите исходные $V_0$ и $S_0$};
\node[actionbox,below=of base] (volume) {Вычтите объём выреза, прибавьте объём добавленной части};
\node[actionbox,below=of volume] (faces) {Выделите изменённые площадки; при необходимости постройте сечение};
\node[choice,below=of faces] (check) {Учтены все исчезнувшие и новые площадки?};
\node[actionbox,right=12mm of check,text width=31mm] (again) {Проверьте торцы, стенки и контакты};
\node[actionbox,below=of check] (sum) {Вычислите $S=S_0-S_-+S_+$};
\node[actionbox,below=of sum] (units) {Проверьте единицы и запишите оба ответа};
\node[start,below=of units] (end) {Конец};
\draw[arr] (start)--(base);\draw[arr] (base)--(volume);\draw[arr] (volume)--(faces);\draw[arr] (faces)--(check);
\draw[arr] (check)--node[above,font=\small]{нет}(again);\draw[arr] (again.north) |- (faces.east);
\draw[arr] (check)--node[left,font=\small]{да}(sum);\draw[arr] (sum)--(units);\draw[arr] (units)--(end);
\end{tikzpicture}
\end{center}
\clearpage
\section{Конус: основные формулы}
Рассматриваем прямой круговой конус. $R$ --- радиус основания, $H$ --- высота, $L$ --- образующая (отрезок от вершины до точки окружности основания).
\begin{center}
\begin{tikzpicture}[scale=.85]
\tkzDefPoint(0,0){A}\tkzDefPoint(6,0){C}\tkzDefPoint(3,0){O}\tkzDefPoint(3,4){B}
\tkzDrawPolygon[color=darkaccent](A,B,C)
\tkzDrawSegments[color=accent](B,O O,C)
\tkzMarkRightAngle[size=.27](B,O,C)
\tkzLabelPoints[below](A,O,C)\tkzLabelPoints[above](B)
\node[left] at (3,2){$H$};\node[above right] at (4.5,2){$L$};\node[below] at (4.5,-.3){$R$};
\end{tikzpicture}
\par\small Осевое сечение конуса.
\end{center}
\[
L=\sqrt{R^2+H^2},\quad \Sos=\pi R^2,\quad V=\frac13\pi R^2H.
\]
\[
\Sb=\pi RL,\qquad \Sp=\pi RL+\pi R^2=\pi R(L+R).
\]
Осевое сечение --- равнобедренный треугольник с основанием $2R$ и высотой $H$. Его площадь:
\[
S_{\text{осев}}=\frac12\cdot2R\cdot H=RH.
\]
\sub{Пример 6. Радиус равен $3$, высота равна $4$}
\[
L=\sqrt{3^2+4^2}=5,\qquad V=\frac13\pi\cdot9\cdot4=12\pi.
\]
\[
\Sb=15\pi,\quad \Sp=24\pi,\quad S_{\text{осев}}=12.
\]
\textbf{Проверяйте:} в формуле объёма используется высота $H$, в формуле боковой площади --- образующая $L$. Диаметр основания равен $2R$.
\clearpage
\section{Сечение конуса, параллельное основанию}
Плоскость, параллельная основанию, отсекает от вершины меньший конус, подобный исходному. Обозначим его высоту через $h$, радиус через $r$, образующую через $l$.
\begin{center}
\begin{tikzpicture}[scale=.85]
\tkzDefPoint(0,0){A}\tkzDefPoint(6,0){C}\tkzDefPoint(3,0){O}\tkzDefPoint(3,5){B}\tkzDefPoint(3,2.5){M}\tkzDefPoint(4.5,2.5){K}\tkzDefPoint(1.5,2.5){J}
\tkzDrawPolygon[color=darkaccent](A,B,C)
\tkzDrawSegments[color=accent,line width=1pt](J,K B,M)
\tkzDrawSegments[dashed](M,O)
\tkzMarkRightAngle[size=.22](B,O,C)\tkzMarkRightAngle[size=.22](B,M,K)
\tkzLabelPoints[above](B)\tkzLabelPoints[below](A,O,C)\tkzLabelPoints[left](J,M)\tkzLabelPoints[right](K)
\node[above] at (3.75,2.5){$r$};\node[below] at (4.5,-.3){$R$};
\draw[<->,accent] (6.5,0)--(6.5,5);\node[right] at (6.5,2.5){$H$};
\draw[<->,accent] (-.5,2.5)--(-.5,5);\node[left] at (-.5,3.75){$h$};
\end{tikzpicture}
\end{center}
Из подобия прямоугольных треугольников $BOC$ и $BMK$:
\[
k=\frac{H}{h}=\frac{R}{r}=\frac{L}{l}>1.
\]
Следовательно,
\[
\frac{V_{\text{бол}}}{V_{\text{мал}}}=\frac{R^2H}{r^2h}=k^3,
\qquad \frac{S_{\text{бок, бол}}}{S_{\text{бок, мал}}}=\frac{RL}{rl}=k^2.
\]
\sub{Пример 7. Объём исходного конуса равен $12$}
Сечение делит высоту пополам. Тогда $h=H/2$, $k=2$ и
\[
V_{\text{мал}}=\frac{12}{2^3}=\frac32=1{,}5.
\]
Объём оставшейся нижней части: $12-1{,}5=10{,}5$. Если требуется площадь боковой поверхности нижней части, вычтите боковую площадь малого конуса из боковой площади большого. Для полной поверхности дополнительно учитываются оба круглых основания нижней части:\[
S_{\text{полн, низ}}=\pi RL-\pi rl+\pi R^2+\pi r^2.
\]
\clearpage
\section*{Тренировочный блок}
\smallskip
10 задач для самостоятельного повторения по желанию. Все размеры даны в одинаковых единицах длины. Записывайте формулы, промежуточные вычисления и ответы.
\sub{Базовый уровень}
\begin{enumerate}[itemsep=5pt]
\item Радиус конуса $5$, высота $12$. Найдите образующую, объём, полную площадь поверхности и площадь осевого сечения.
\item Диагональ основания правильной четырёхугольной пирамиды равна $12$, боковое ребро равно $10$. Найдите объём пирамиды.
\item Из угла параллелепипеда $8\times5\times4$ вырезали куб с ребром $2$, имеющий три грани на исходной внешней поверхности. Найдите объём и площадь поверхности остатка.
\item Ребро правильного тетраэдра увеличили в четыре раза. Во сколько раз увеличились полная площадь поверхности и объём?
\item Сторону квадратного основания пирамиды увеличили в четыре раза, высоту --- в два раза. Во сколько раз увеличился объём?
\end{enumerate}
\sub{Уровень выше среднего}
\begin{enumerate}[start=6,itemsep=5pt]
\item Сторона основания правильной четырёхугольной пирамиды равна $10$, угол между боковой гранью и основанием --- $45^\circ$. Найдите объём и полную площадь поверхности.
\item Объём конуса равен $81$. Сечение, параллельное основанию, отсекает от вершины конус высотой в треть исходной. Найдите объёмы обеих частей.
\item Радиусы оснований большого и отсечённого конусов равны $9$ и $3$. Боковая площадь большого конуса равна $108\pi$. Найдите боковую площадь малого конуса и нижней части. Сечение параллельно основанию.
\item В параллелепипеде $9\times6\times5$ вырезали сквозной паз $9\times2\times3$, открытый сверху и на двух торцах, строго между боковыми гранями. Найдите объём и полную площадь поверхности остатка.
\end{enumerate}
\sub{Повышенный уровень}
\begin{enumerate}[start=10]
\item В параллелепипеде $7\times5\times4$ вырезали сквозной паз $7\times2\times2$, открытый сверху и на двух торцах. По ширине верхней грани оставлены полосы шириной $2$ и $1$. На широкой полосе закрепили куб с ребром $2$, полностью прилегающий нижней гранью к верху, без перекрытия паза. Найдите объём и полную площадь поверхности фигуры.
\end{enumerate}
\clearpage
\section*{Ответы}
\renewcommand{\arraystretch}{1.22}
\begin{tabularx}{\linewidth}{@{}cX@{}}
\toprule
№ & Ответ\\\midrule
1 & $L=13$, $V=100\pi$, $\Sp=90\pi$, $S_{\text{осев}}=60$.\\
2 & $V=192$.\\
3 & $V=152$, $\Sp=184$.\\
4 & Площадь в $16$ раз; объём в $64$ раза.\\
5 & В $32$ раза.\\
6 & $V=\dfrac{500}{3}$, $\Sp=100(1+\sqrt2)$.\\
7 & Малый конус: $3$; нижняя часть: $78$.\\
8 & Малый конус: $12\pi$; нижняя часть: $96\pi$.\\
9 & $V=216$, $\Sp=300$.\\
10 & $V=120$, $\Sp=202$.\\\bottomrule
\end{tabularx}
\vspace{8mm}
\sub{Итоговая памятка}
\begin{tabularx}{\linewidth}{@{}p{.38\linewidth}X@{}}
\toprule
Сигнал в задаче & Полезное действие\\\midrule
Диагональ квадрата & $\Sos=d^2/2$, $OA=d/2$.\\
Угол грани с основанием & Постройте $\triangle SOM$; используйте $\tg\alpha$ и $\cos\alpha$.\\
Подобные тела & Линейный масштаб $k$, площади $k^2$, объёмы $k^3$.\\
Угловой вырез & Вычтите объём, сравните три пары площадок.\\
Сквозной паз & Проверьте торцы и продольные стенки.\\
Присоединённый куб & Исключите обе контактные поверхности.\\
Параллельное сечение конуса & Определите $k=H/h$ и нужную степень.\\\bottomrule
\end{tabularx}
\end{document}
